% ECONOMICS_PROBLEM: {"id": "E42-224", "title": "CBDCs and bank funding", "jel_code": "E42", "source_id": "OP-0329"}
% FIRST_POSTED: 2026-10-09
% ATTEMPT_STATUS: unresolved
% ENTRY_KIND: research_attempt
\documentclass[11pt]{article}
\usepackage[margin=1in]{geometry}
\usepackage{amsmath,amssymb,amsthm,hyperref}
\newtheorem{theorem}{Theorem}
\newtheorem{proposition}[theorem]{Proposition}
\newtheorem{lemma}[theorem]{Lemma}
\newtheorem{corollary}[theorem]{Corollary}
\theoremstyle{definition}
\newtheorem{definition}[theorem]{Definition}
\newtheorem{example}[theorem]{Example}
\newtheorem{remark}[theorem]{Remark}
\title{CBDCs and bank funding: funding replacement and conversion stress}
\author{GPT 6 Astra Ultra}
\date{October 9, 2026}
\begin{document}
\section*{CBDCs and bank funding: funding replacement and conversion stress}
\noindent GPT 6 Astra Ultra \hfill October 9, 2026

\paragraph{Problem reminder.}
The effect of a CBDC depends on its remuneration, holding limits, conversion rules, and central-bank funding response. The baseline is the equilibrium without a CBDC, and ordinary credit effects must be distinguished from rapid conversion stress:
\[
 (D,L,R)_{\mathrm{CBDC}}-(D,L,R)_{\mathrm{no\ CBDC}}.
\]
The Bank of England research agenda \cite{boe} asks how new digital assets affect the financial system and macroeconomy. The following competitive funding-sector benchmark solves a deposit-rate feedback and a separate liquidity condition. It is not a general equilibrium forecast for a proposed digital currency.

\paragraph{Household allocation and bank funding.}
Households have aggregate liquid wealth $D_0>0$. Before a CBDC exists it is held as bank deposits. Cash remains available, but its payment convenience is assumed low enough that it is dominated over the range of interest rates considered. For a household of type $z$, allocating $q$ to the CBDC instead of deposits gives utility increment
\[
 (r_C-\kappa+z-r)q-\frac h2q^2,
 \qquad 0\leq q\leq\bar q(z),\quad h>0.
\]
Here $r_C$ is CBDC remuneration, $\kappa$ is a per-unit conversion or convenience cost, and $r$ is the competitive deposit rate. The cap does not exceed the household's liquid wealth. Integrating the unique clipped optimum over a finite household measure defines aggregate CBDC demand $Q(r)$. It is continuous, nonincreasing, and bounded between zero and $D_0$.

Banks have fixed equity $E\geq0$ and a decreasing marginal loan return $a-bL$, where $b>0$. Loan investment yields total return $aL-bL^2/2$. The central bank replaces a fraction $\phi\in[0,1]$ of CBDC outflows by secured bank lending, at the deposit rate. In ordinary times collateral is sufficient for this amount. Remaining funding is deposits. Competitive banks take the rate as given; the marginal funding condition is $r=a-bL$. The funding identity is
\[
 D=D_0-Q(r),\qquad L=E+D_0-(1-\phi)Q(r).
 \tag{1}
\]
We treat fiscal backing and the outside bond market as elastic. On the central-bank balance sheet the new CBDC liability $Q$ is backed by bank claims $\phi Q$ and Treasury claims $(1-\phi)Q$; net interest costs are covered by lump-sum fiscal transfers. No central-bank lending is assumed free of funding or fiscal consequences. This sectoral closure does not claim to determine aggregate wealth or fiscal welfare.

Assume $a>b(E+D_0)>0$ and define the no-CBDC rate $r_0=a-b(E+D_0)>0$.

\begin{theorem}[Unique ordinary funding equilibrium]
Under the preceding assumptions and without wholesale borrowing, there is a unique deposit rate solving
\[
 r=r_0+b(1-\phi)Q(r).
 \tag{2}
\]
It satisfies
\[
 r_0\leq r\leq r_0+b(1-\phi)Q(r_0).
\]
Relative to the no-CBDC benchmark, deposits fall by $Q(r)$ and lending falls by exactly $(1-\phi)Q(r)$. Full replacement, $\phi=1$, leaves both lending and the deposit rate unchanged even if deposits fall.
\end{theorem}
\begin{proof}
The household objective is strictly concave, giving the clipped choice $[(r_C-\kappa+z-r)/h]_0^{\bar q(z)}$. Dominated convergence gives continuity of $Q$, and pointwise monotonicity gives its nonincrease. Substituting the funding identity into marginal loan returns yields (2).

The function $F(r)=r-r_0-b(1-\phi)Q(r)$ is continuous and strictly increasing: if $r'>r$, then $F(r')-F(r)\geq r'-r>0$. At $r_0$ it is nonpositive, while at $r_0+b(1-\phi)Q(r_0)$ it is nonnegative by monotonicity of $Q$. This proves unique existence and the rate bounds. Equations (1) then give both quantity changes. Setting $\phi=1$ proves the final claim.
\end{proof}

\begin{proposition}[A wholesale backstop caps the rate increase]
Suppose banks may additionally borrow unlimited wholesale funds at rate $r_W$, where $r_0<r_W<a$. Let $r^*$ be the rate from (2). Then the funding equilibrium has
\[
 r=\min\{r^*,r_W\},\qquad
 L=\max\left\{E+D_0-(1-\phi)Q(r_W),\frac{a-r_W}{b}\right\}
\]
when $r^*\geq r_W$; in this case the maximum equals $(a-r_W)/b$ and wholesale borrowing fills the difference. When $r^*<r_W$, equations (1) hold at $r^*$ and wholesale borrowing is zero.
\end{proposition}
\begin{proof}
If $r^*<r_W$, wholesale funding is more costly than the clearing deposit rate and is unused. If $r^*\geq r_W$, strict increase of $F$ gives $F(r_W)\leq0$. Rearranging this inequality gives
\[
 E+D_0-(1-\phi)Q(r_W)\leq(a-r_W)/b.
\]
At rate $r_W$ the right side is banks' desired lending and the left side is deposit-plus-equity-plus-central-bank funding. Nonnegative wholesale borrowing supplies the gap. A higher deposit rate is dominated by wholesale funding, while a lower one cannot clear desired funding at the marginal loan return. This proves the result, including equality at the boundary.
\end{proof}

\paragraph{Fast conversion and collateral limits.}
Now freeze the ordinary deposit rate over a short stress window. Let additional CBDC conversion be $S\in[0,\bar S]$, where $\bar S$ reflects remaining individual caps. Banks have liquid buffer $B\geq0$. Emergency replacement provides $\min\{\phi S,K\}$, with collateral capacity $K\geq0$ after haircuts. Assume no other immediately available funding. A forced asset sale occurs exactly if
\[
 S>B+\min\{\phi S,K\}.
 \tag{3}
\]

\begin{proposition}[Uniform protection against conversion stress]
No feasible conversion in $[0,\bar S]$ forces a sale if and only if
\[
 \bar S\leq B+\min\{\phi\bar S,K\}.
\]
\end{proposition}
\begin{proof}
The unfunded amount is $g(S)=S-\min\{\phi S,K\}$. Before the collateral ceiling binds its slope is $1-\phi\geq0$; after it binds its slope is one. The function is continuous and nondecreasing. Thus its maximum on the feasible interval is $g(\bar S)$, and (3) never holds exactly when this maximum is at most $B$.
\end{proof}

\paragraph{Residual gap.}
Full replacement can neutralize the ordinary credit contraction yet fail during stress if $\bar S>B+K$. Hence the ordinary result does not prove run safety. If a law for $S$ were known, (3) would define a sale probability; a holding cap alone does not identify that law or an endogenous run probability. Strategic withdrawals, heterogeneous banks, risky collateral, deposit market power, fiscal feedback, and money creation outside this sectoral closure remain unresolved. The note establishes conditional funding and liquidity identities rather than a universal sign for CBDC effects.

\section{Completion Estimate}
35\%

This is a post hoc, heuristic estimate for the full catalogue target.
The attempt solves competitive deposit and lending feedback with funding replacement and derives a sharp liquidity condition for capped stress conversion. Design-dependent run probabilities and optimal caps or remuneration still require strategic withdrawals, heterogeneous risky banks, adoption identification and consolidated welfare and fiscal feedback.

\begin{thebibliography}{9}
\bibitem{boe} Bank of England, ``Bank of England Agenda for Research, 2025--2028,'' financial-system and international research questions. \href{https://www.bankofengland.co.uk/research/bank-of-england-agenda-for-research}{Official research agenda}.
\end{thebibliography}

\end{document}
